\emph{Basa tataran kapisan} dumadi saka \emph{simbol konstanta},
\emph{simbol fungsi}, lan \emph{simbol predikat}. Simbol fungsi
lan predikat nduwèni cacah argumèn kang wis ditemtokaké.
Simbol konstanta ora mbutuhaké argumèn. Upamané, ing
\emph{basa aritmetika}, ana siji simbol konstanta~$\Obj 0$,
siji simbol fungsi siji-panggonan $\prime$, rong simbol fungsi
rong-panggonan $+$ lan~$\times$, lan siji simbol predikat
rong-panggonan~$<$. Saka \emph{variabèl}, simbol konstanta,
lan simbol fungsi, kita mbentuk \emph{terma}. Saka terma,
simbol predikat, lan \emph{simbol identitas}~$\eq$, kita
mbentuk \emph{formula atomik}. Saka formula atomik, kita
mbentuk formula nganggo panggandhèng logis $\lnot$, $\lor$,
$\land$, $\lif$, $\liff$ lan kuantor $\lforall$ lan $\lexists$.
Amarga tandha kurung kang dibutuhaké tansah dienggo, saben formula
mung nduwèni siji cara maca. Mula kita bisa mènèhi definisi
nganggo induksi tumrap struktur formula.

Kedadéan variabèl ing formula kadhang dikuwasani kuantor kang
cocog. Kedadéan iku \emph{kaiket} yèn ana ing \emph{lingkup}
kuantor kang cocog, lan \emph{bebas} yèn ora. Konsèp-konsèp
iki ditetepaké kanthi induksi. Operasi \emph{substitusi} sawijining
terma kanggo variabèl ing formula uga ditetepaké kanthi induksi.
Formula tanpa kedadéan variabèl bebas diarani \emph{ukara}.
